\section*{Chapter \ref{ch:g8:dictionaries} --- Dictionaries}

\begin{solution}{exo:g8:dictionaries:1}
\code{d["b"]} is 2. \code{d["d"]} stops the \omterm{def:g4:programs-and-operating-system:program}{program} with a \code{KeyError}:
there is no \omterm{def:g8:dictionaries:dictionary}{key} \code{"d"}.
\end{solution}

\begin{solution}{exo:g8:dictionaries:2}
The second \omterm{def:g6:variables-and-input:assignment}{assignment} changes the \omterm{def:g6:variables-and-input:value}{value} of the \omterm{def:g8:dictionaries:dictionary}{key} \code{"x"}, which is
already there: \code{len(d)} is 1 and \code{d["x"]} is 5.
\end{solution}

\begin{solution}{exo:g8:dictionaries:3}
\code{Sami}, then \code{Leo}: the two names whose age is 12, in the order of
the \omterm{def:g8:dictionaries:dictionary}{dictionary}.
\end{solution}

\begin{solution}{exo:g8:dictionaries:4}
\omcode{code/grade-8/02-dictionaries/exercises/capitals.py}{1}{8}{A lookup with
a default.}\label{lst:g8:dictionaries:sol-capitals}

It prints \code{Tokyo / not known}.
\end{solution}

\begin{solution}{exo:g8:dictionaries:5}
\code{french[word]} needs every word to be a \omterm{def:g8:dictionaries:dictionary}{key}. With
\code{french.get(word, word)}, a missing word gives itself back:

\omcode{code/grade-8/02-dictionaries/exercises/translate_fixed.py}{4}{6}{The
\omterm{def:g4:algorithms-repeat-and-decide:loop}{loop}, fixed.}\label{lst:g8:dictionaries:sol-translate}

It prints \code{le chat mange le mouse}.
\end{solution}

\begin{solution}{exo:g8:dictionaries:6}
\code{counts[word] = 1} sets the count to 1 each time, so a word seen twice
still counts 1. The counting pattern adds one to the old count:
\code{counts[word] = counts.get(word, 0) + 1}, which gives 2 for \code{a} and
for \code{rose}.
\end{solution}

\begin{solution}{exo:g8:dictionaries:7}
Count the words, then keep the best so far:

\omcode{code/grade-8/02-dictionaries/exercises/most_frequent.py}{1}{12}{The most
frequent word.}\label{lst:g8:dictionaries:sol-frequent}

It prints \code{the}, which appears four times.
\end{solution}

\begin{solution}{exo:g8:dictionaries:8}
\omcode{code/grade-8/02-dictionaries/exercises/invert.py}{1}{9}{\omterm{def:g6:variables-and-input:value}{Values} become
\omterm{def:g8:dictionaries:dictionary}{keys}.}\label{lst:g8:dictionaries:sol-invert}

The \omterm{def:g6:variables-and-input:value}{values} become \omterm{def:g8:dictionaries:dictionary}{keys}, and \omterm{def:g8:dictionaries:dictionary}{keys} are all different: if two \omterm{def:g8:dictionaries:dictionary}{keys} had the same
\omterm{def:g6:variables-and-input:value}{value}, the second \omterm{def:g8:dictionaries:dictionary}{entry} would replace the first in the result.
\end{solution}

\begin{solution}{exo:g8:dictionaries:9}
Three entries: \code{a} 3, \code{rose} 3 and \code{is} 2.
\end{solution}

\begin{solution}{exo:g8:dictionaries:10}
\omcode{code/grade-8/02-dictionaries/exercises/oldest.py}{8}{15}{The oldest
book.}\label{lst:g8:dictionaries:sol-oldest}

It prints \code{The Lighthouse Cat} (1998).
\end{solution}

\begin{solution}{exo:g8:dictionaries:11}
The first method walks the $b$ ballots once per candidate: $c \times b$ names.
The second walks them once: $b$ names. For 1\,000 ballots and 50 candidates,
50\,000 against 1\,000: the \omterm{def:g8:dictionaries:dictionary}{dictionary} looks at 50 times fewer names, as
many times fewer as there are candidates.
\end{solution}

\begin{solution}{exo:g8:dictionaries:12}
A word whose first letter is new starts an empty \omterm{def:g7:lists:list}{list} first:

\omcode{code/grade-8/02-dictionaries/exercises/group.py}{1}{11}{A \omterm{def:g8:dictionaries:dictionary}{dictionary} of
\omterm{def:g7:lists:list}{lists}.}\label{lst:g8:dictionaries:sol-group}

It prints \code{{'a': ['apple', 'avocado'], 'b': ['bean', 'banana'], 'c': ['cherry']}}.
\end{solution}

\begin{solution}{pb:g8:dictionaries:1}
\textbf{1.} Chloe 4, Ada 2, Eva 2, Dev 2, Ben 1, Femi 1: Chloe leads.

\textbf{2.} One tally line per name, added the first time the name appears:
exactly what a \omterm{def:g8:dictionaries:dictionary}{dictionary} of counts does.

\textbf{3.} Every name looked at; the inner \omterm{def:g4:algorithms-repeat-and-decide:loop}{loop} runs once per ballot, for
each of the six candidates.

\textbf{4.} Chloe wins:

\omoutput{code/grade-8/02-dictionaries/outputs/votes_list.txt}{Six walks of 40
ballots.}\label{out:g8:dictionaries:votes-list}

\textbf{5.} \code{counts[ballot] = counts.get(ballot, 0) + 1}.

\textbf{6.} It walks the counts again and keeps the largest count of any
name other than the winner; the margin is the difference.

\textbf{7.} They agree, with 40 names looked at instead of 240:

\omoutput{code/grade-8/02-dictionaries/outputs/votes_dict.txt}{One walk of 40
ballots.}\label{out:g8:dictionaries:votes-dict}

\textbf{8.} A \omterm{def:g8:dictionaries:dictionary}{key} enters the \omterm{def:g8:dictionaries:dictionary}{dictionary} when its name first appears on a
ballot: Chloe is on the first ballot, Ada on the second.

\textbf{9.} $6 \times 40 = 240$ and 40; $20 \times 400 = 8\,000$ and 400.

\textbf{10.}

\omoutput{code/grade-8/02-dictionaries/outputs/votes_compare.txt}{The two
methods on two elections.}\label{out:g8:dictionaries:compare}

The \omterm{def:g7:lists:list}{list} method looks at $c \times b$ names, the \omterm{def:g8:dictionaries:dictionary}{dictionary} at $b$.

\textbf{11.} 480 million names with \omterm{def:g7:lists:list}{lists}, 40 million with a \omterm{def:g8:dictionaries:dictionary}{dictionary}.

\textbf{12.} \textbf{Result.} Counting with one walk per candidate looks at
candidates $\times$ ballots names (240 for the class, 8\,000 for the larger
election); counting with a \omterm{def:g8:dictionaries:dictionary}{dictionary} looks at each ballot once (40, then
400).
\end{solution}
